\documentclass[11pt,a4paper]{article}
\usepackage[margin=22mm,headheight=15pt]{geometry}
\usepackage[T1]{fontenc}
\usepackage{lmodern,microtype,array,booktabs,longtable,enumitem,fancyhdr,listings,xcolor}
\usepackage[hidelinks]{hyperref}
\input{build-info.tex}
\definecolor{navy}{HTML}{17324D}
\definecolor{pale}{HTML}{EDF3F8}
\pagestyle{fancy}\fancyhf{}
\fancyhead[L]{\small\textbf{COS341 2026 / Group 2 / P2}}
\fancyhead[R]{\small First Class}
\fancyfoot[L]{\footnotesize Group 2 user manual}
\fancyfoot[R]{\thepage}
\setlength{\emergencystretch}{3em}
\setlength{\parindent}{0pt}\setlength{\parskip}{6pt}
\setlist[itemize]{leftmargin=*,itemsep=2pt,topsep=3pt}
\lstset{basicstyle=\ttfamily\small,breaklines=true,columns=fullflexible,keepspaces=true,frame=single,rulecolor=\color{navy},backgroundcolor=\color{pale},aboveskip=5pt,belowskip=5pt}
\newcommand{\heading}[1]{\vspace{5pt}{\large\bfseries\color{navy}#1}\par}
\newcommand{\signrow}[2]{\par\noindent\textbf{#1}\quad\texttt{#2}\par\vspace{14pt}Signature: \rule{92mm}{0.4pt}\quad Date: \rule{25mm}{0.4pt}\par\vspace{18pt}}
\begin{document}
{\LARGE\bfseries\color{navy} Group 2 SPL compiler}\par
{\large COS341 2026 --- Practical 2 user manual}\par
I have put this manual together to explain how I run our compiler and the C++ program it produces. Our practical includes Parts 2a, 2b, 3a and 3b, which cover name and scope analysis, type checking, code generation and function inlining. I have prepared it for the \textbf{First Class category}.

\heading{1. Choose the executable for your computer}
I first extract the ZIP and open a terminal \textbf{inside the \texttt{group-2} folder}. I then choose the executable for my operating system from the table below. The compiler includes its Python runtime. I use a separate C++ compiler to build and run the program it generates.

\begin{tabular}{@{}p{25mm}p{50mm}p{69mm}@{}}\toprule
Computer & Compiler file & Supported system\\\midrule
Windows x64 & \texttt{group-2.exe} & Windows x64; tested on Windows Server 2022\\
Linux x64 & \texttt{group-2} & Linux x86-64; Ubuntu 22.04 baseline (glibc 2.35+)\\
macOS Apple Silicon & \texttt{macos/group-2} & macOS arm64; macOS 14 or later\\\bottomrule
\end{tabular}\par\medskip
The Linux and macOS files have \textbf{no extension}. I use the file for the correct operating system, even if Finder describes both as a ``Unix Executable File''. The macOS build is signed locally rather than notarized by Apple, so macOS may ask me to allow it to open.

\heading{2. Windows (PowerShell or Command Prompt)}
\begin{lstlisting}
.\group-2.exe
g++ -std=c++17 program.cpp -o program.exe
.\program.exe
\end{lstlisting}
\heading{3. Linux}
\begin{lstlisting}
chmod +x group-2
./group-2
g++ -std=c++17 program.cpp -o program
./program
\end{lstlisting}
\heading{4. macOS Apple Silicon}
\begin{lstlisting}
chmod +x macos/group-2
./macos/group-2
g++ -std=c++17 program.cpp -o program
./program
\end{lstlisting}
\textbf{Expected example output:} \texttt{result:3:4}, without an extra newline. I have included \texttt{SPL.txt} as a working example. It calls a numeric function twice and uses a nested procedure to update its parent's local variable inside a loop.

\newpage
\heading{5. Input, output and errors}
When I run the compiler without an argument, it reads \texttt{SPL.txt} from the current folder. To compile a different ASCII source file, I pass its path as the only argument. I put quotation marks around a path that contains spaces:
\begin{lstlisting}
.\group-2.exe "path\to\input with spaces.spl"
./group-2 "path/to/input with spaces.spl"
./macos/group-2 "path/to/input with spaces.spl"
\end{lstlisting}
If compilation succeeds, the compiler returns status 0 without printing a message. I get two files in the current folder: \texttt{tree.xml}, containing the syntax tree, and \texttt{program.cpp}, containing the generated C++ code. The XML keeps the original tokens and node IDs. The C++ file includes the function code mine as comments and the executable code in one \texttt{main}. I compile and run this C++ file to see my SPL program's output.

\textbf{Writing the input.} I separate every token, including the last one, with a space or newline (LF, CR or CRLF). The language accepts ASCII input, lowercase names beginning with \texttt{\#}, and strings without spaces. Tabs and a literal \texttt{\$} are invalid. I write zero as \texttt{0}; decimals must end in a nonzero fractional digit. I keep the final newline when editing the example.

\begin{tabular}{@{}p{17mm}p{128mm}@{}}\toprule
Status & Meaning\\\midrule
0 & Complete compiler success; fresh XML and C++ outputs are available.\\
2 / 3 & Incorrect command usage / input or path-alias problem.\\
4 / 5 / 6 & Lexical / syntax / tree error.\\
7 / 8 & Name or scope error / type or whole-program numeric conflict.\\
9 & Code-generation or output-publication error.\\
70 & Unexpected internal failure or an OS error preventing output cleanup.\\\bottomrule
\end{tabular}\par\medskip
If there is an error, I get a message on stderr showing the input file and the relevant location. A failed full compilation normally removes both output files. The compiler checks for input/output path collisions before deleting anything, so it preserves my source file. It also reports any operating-system error that prevents it from clearing an old output.

\heading{6. Trying an invalid program}
I have included a second example so I can check how the compiler reports a name error:
\begin{lstlisting}
.\group-2.exe examples\invalid-names.spl
./group-2 examples/invalid-names.spl
\end{lstlisting}
On macOS I use \texttt{./macos/group-2} in the same command. This example uses \texttt{\#missing} without declaring it, so I expect status 7 and \texttt{NAME-UNDECLARED-VARIABLE}. Neither output file should remain. I can run the compiler again without an argument to restore the valid example's outputs.

\heading{7. Compiling the generated C++}
\textbf{I recommend \texttt{g++} with C++17 for assessment.} Our output uses \texttt{std::cout}, and Announcement 34 allows us to describe it as C++. The generated \texttt{program.cpp} can be inspected before running it. Our tests cover both GNU \texttt{g++} 13 and \texttt{gcc} 13. If I use \texttt{gcc}, I select C++ and link its standard library explicitly:
\begin{lstlisting}
gcc -x c++ -std=c++17 program.cpp -lstdc++ -o program
\end{lstlisting}
On Windows I use \texttt{-o program.exe}. My GNU toolchain needs the C++ frontend and \texttt{libstdc++}. I only need this toolchain to build the generated program; our SPL compiler itself runs separately.

\newpage
\heading{8. Files I have included}
The ZIP contains the following files:
\begin{lstlisting}
group-2/
  group-2.pdf              This manual and two personal declarations
  group-2.exe              Windows x64 compiler
  group-2                  Linux x64 compiler (no extension)
  macos/group-2            macOS arm64 compiler (no extension)
  SPL.txt                  Valid repeated/nested-call example
  examples/invalid-names.spl
  manifest.json            Source commit, platforms and SHA-256 hashes
\end{lstlisting}
I have included a SHA-256 checksum alongside the ZIP so I can check that the file has not changed. The \texttt{manifest.json} file records the build version and checksums of the files inside. I keep the Windows, Linux and macOS executables under their supplied names and use the correct one for the computer running the assessment.

\heading{9. How our compiler handles SPL}
I have summarised the main rules below to explain the compiler's behaviour:
\begin{itemize}
\item A variable refers to its nearest declaration in the current scope or an enclosing scope. It cannot refer sideways to a declaration in another function. A function call must refer to a function declared directly inside the current P. Duplicate declarations, a local masking its own parameter, overloading and recursion are rejected.
\item Functions are inlined. Arguments are evaluated from left to right and saved before being copied into the function's parameters. The number of arguments must match the number of parameters. A numeric call returns its result to the caller; a void call has no numeric result. Boolean \texttt{and}/\texttt{or} use short circuiting, and the generated control flow uses comparisons and jumps.
\item The whole program uses \texttt{float} if it contains division or a decimal number; otherwise it uses \texttt{int}. A program containing both \texttt{mod} and \texttt{div}, or \texttt{mod} and a decimal number, is rejected. Text inside strings does not affect this check.
\item The output uses native integers and single-precision floats, with the normal \texttt{std::cout} format (six significant digits for floats). In my examples, I initialise variables, keep values within the native number ranges and avoid division by zero. Inlining is limited to 1,000,000 generated instructions.
\item Variables and functions have separate naming tables. A return expression uses the function body's scope. Scope and type information are kept internally, while the output files are \texttt{tree.xml} and \texttt{program.cpp}.
\end{itemize}
\heading{10. Preparing our submission}
Before I submit, each member needs to review the files and personally sign \emph{both} declarations on the following pages. I will include the signed pages in this manual, check \texttt{group-2.zip}, and upload it \textbf{only to the First Class slot} on 25 October 2026 before the displayed cutoff. I will keep the receipt and compare the downloaded ZIP's checksum with my original copy.
\par\medskip
{\footnotesize Build reference: \texttt{\SourceCommit}}\par

\newpage
{\LARGE\bfseries\color{navy} Declaration 1: submission category}\par
\textbf{COS341 2026 --- Group 2 --- Part 2}\par
By signing below, I confirm that I agree to submit our Group 2 practical in the \textbf{First Class category}, for a maximum of \textbf{6/6 Part 2 points} in addition to our parser points. I understand that our submission includes Parts \textbf{2a, 2b, 3a and 3b}, including the generation of readable C++ code that can be compiled and run.

Each member signs and dates their own row.
\par\vspace{14pt}
\signrow{Rudolph}{u20598425}
\signrow{Kuda}{u20620111}
\signrow{Cherise}{u21434019}
\signrow{Michele}{u21448842}
\vfill
\textit{This is our declaration of submission category, as required by Announcement 33.}

\newpage
{\LARGE\bfseries\color{navy} Declaration 2: complete ZIP contents}\par
\textbf{COS341 2026 --- Group 2 --- Part 2}\par
By signing below, I confirm that I am fully aware of the \textbf{whole contents of our final submission ZIP}. This includes the compiler executables, this manual and its declarations, the example programs and \texttt{manifest.json}. I have reviewed these files and \textbf{agree to their submission in this form} in the First Class category.

Each member signs and dates their own row after reviewing the same final ZIP. If its contents change, we will review the changes before submitting.
\par\vspace{14pt}
\signrow{Rudolph}{u20598425}
\signrow{Kuda}{u20620111}
\signrow{Cherise}{u21434019}
\signrow{Michele}{u21448842}
\vfill
\textit{This declaration records my agreement to the ZIP contents under Announcement 33.}
\end{document}
